%% Section 4 -- Reconstruction & Reconciliation Algorithms
%% Ported from docs/DRAFT_PAPER.md Section 4. The algorithm body follows
%% docs/PAPER_BLUEPRINT.md 5 verbatim.
%%
%% LOCKED: the procedure is GroupIntoReconciliationClusters. The retired name
%% PartitionIntoEquivalenceClasses must never reappear here -- ~_dedup is a
%% similarity relation, not an equivalence relation, and it yields clusters,
%% not classes.
%%
%% LOCKED (2026-08-26, formal-consistency pass): the in-cluster predicate is
%% HasUnresolvedConflict, not the retired IsAmbiguous -- the question is what
%% remains unsettled after the declared rules run, not what was ambiguous at
%% some point during processing. DetectCrossClusterConflicts must stay a line
%% of Algorithm 1 and must stay AFTER the loop: it compares canonical records,
%% so it cannot run until every cluster has produced one. Line numbering in the
%% walkthrough prose below depends on this; renumber it if the block changes.

\section{Reconstruction \& Reconciliation Algorithms}
\label{sec:algorithms}

\subsection{Ordering and timestamp granularity}
\label{sec:algorithms:ordering}

Two distinct timestamp tolerances operate in Synthesis, and conflating them is
the easiest way to misread the specification.

The first is the \textbf{deduplication window} $\Delta t_{\text{granularity}}$,
which governs whether two records may describe the same message. Exports quantise
time. A platform that writes minute-resolution timestamps into one export and
second-resolution into another will represent one message with two different
values, so an exact timestamp match would fail to group records that plainly
belong together.

$\Delta t_{\text{granularity}}$ is a property of the input, not of the run. It
should be set to the coarsest timestamp quantum among the fragments being
reconciled, which an implementation may infer from the exports or take from the
operator, and it must be recorded in $\Omega_{\text{stats}}$ either way.
Invariant 3 constrains an implementation to reproduce its own output, so the
value has to be a function of the input rather than of when the run happened; a
window that varied between runs on identical fragments would break the invariant
directly, and an implementation that lets an operator supply it must therefore
treat it as part of the input it is being deterministic about.

The second is the \textbf{coarse resolution} $\Delta t_{\text{conflict}}$ used
for conflict detection, a deliberately blunter bucket used to decide which record
pairs are close enough that a contradiction between them is worth reporting. Like
the deduplication window it is a declared parameter rather than a constant, and
its default is one minute. That default is not an empirical finding about export
formats in general, and Section~\ref{sec:implementation:inputs} is explicit that
the corpus is not a field survey; it is chosen to match the timestamp resolution
of the export format the reference implementation targets, which records minutes
and not seconds. An implementation targeting a format with finer or coarser
timestamps should declare a different value, and
Section~\ref{sec:carta:conformance} requires it to say which.

The two tolerances are not independent. A conforming implementation must satisfy

\begin{equation}
  \Delta t_{\text{granularity}} \le \Delta t_{\text{conflict}}
  \label{eq:tolerances}
\end{equation}

and this is a constraint rather than an observation. Were the deduplication
window the wider of the two, records grouped into a single cluster could be
further apart than the bucket used to decide whether two \emph{different}
messages contradict each other, and the same pair of records could be
simultaneously treated as one message and as two conflicting ones. Widening one
therefore requires widening the other.

Ordering comes before either. Records are sorted primarily by $t_{\text{utc}}$
ascending and secondarily by source sequence index. The secondary key is what
makes Invariant 3 attainable: records sharing a timestamp would otherwise be
ordered by whatever the runtime happened to do, and two runs could disagree.
Determinism here is bought by a tie-break rule, not assumed.

\subsection{The similarity criterion}
\label{sec:algorithms:similarity}

Two normalised records are related by $\sim_{\text{dedup}}$ when they agree on
sender, body, and attachment, and their timestamps fall within the deduplication
window:

\begin{equation}
  \begin{split}
    r'_a \sim_{\text{dedup}} r'_b \iff{}
      &|t_{\text{utc}}(r'_a) - t_{\text{utc}}(r'_b)|
        \le \Delta t_{\text{granularity}} \\
      &{}\land\; s(r'_a) = s(r'_b)
       \;\land\; b(r'_a) = b(r'_b)
       \;\land\; a(r'_a) \cong a(r'_b)
  \end{split}
  \label{eq:dedup}
\end{equation}

This is a \textbf{similarity relation} and we do not claim it is an equivalence
relation. Proximity criteria are not generally transitive. Three records spaced
just under $\Delta t_{\text{granularity}}$ apart in sequence will satisfy the
relation pairwise along the chain while the outer two fall outside each other's
window. Treating the relation as an equivalence and speaking of equivalence
classes or a quotient set would assert a partition that the criterion does not
deliver.

Records satisfying the criterion are therefore grouped into
\textbf{reconciliation clusters}, and the specification fixes the criterion while
leaving the clustering procedure open. Transitive closure with a bounded cluster
diameter and connected components over a similarity graph are both admissible,
and they can disagree on the chain case above. Choosing between them is an
implementation decision deferred to the implementation phase, and a conforming
implementation must state which rule it applies, since Invariant 3 constrains a
given implementation to behave the same way twice rather than requiring all
implementations to agree with each other.

For each cluster $C$, one \textbf{canonical} record $u$ is synthesised. Canonical
here means the deterministic representative of $C$ under the implementation's
declared reconciliation policy, and not a globally unique record that every
conforming implementation would agree on. Because the clustering procedure is
open, two conforming implementations may cluster the same input differently and
therefore produce different representatives; Invariant 3 requires each of them
to be self-consistent, not to agree with the other.

Synthesis is field-wise, and it is the second open parameter of the
specification. For each field of the canonical schema, an implementation applies
a fixed precedence rule over the observations in $C$. Where those observations
carry compatible values, the rule selects one deterministically. Where they are
incompatible, the rule still yields a value, and the cluster additionally emits
a conflict descriptor, so that a choice made between contradictory observations
is reported rather than made silently. The specification fixes that shape and
leaves the precedence itself open on the same terms as the clustering procedure:
a conforming implementation must state the rule it applies. Fixing one
precedence here would amount to a claim about which export is the more
trustworthy source, and this paper has no evidence supporting such a claim.

The provenance of $u$ is the union of the provenance of every record in the
cluster:

\begin{equation}
  \mathrm{prov}(u) = \bigcup_{r' \in C} \mathrm{prov}(r')
  \label{eq:provunion}
\end{equation}

The union is what preserves Invariant 2 across deduplication. Collapsing three
observations into one record without unioning their provenance would discard two
of the three origins and leave the archive claiming a single source for something
seen three times.

\subsection{The Synthesis procedure}
\label{sec:algorithms:procedure}

\begin{algorithm}[htbp]
\caption{CARTA Synthesis and Audit Report Construction}
\label{alg:carta_synthesis}
\begin{algorithmic}[1]
\Require Normalized fragment record sets
  $\mathcal{R}' = \{R'_1, R'_2, \dots, R'_n\}$,
  quarantine sets $\mathcal{Q} = \{Q_1, \dots, Q_n\}$,
  deduplication window $\Delta t$
\Ensure Unified archive $\mathcal{U}$, reconstruction report $\Omega$
\State $\mathcal{U} \leftarrow \emptyset$,
       $\mathcal{C}_{\text{conflicts}} \leftarrow \emptyset$,
       $\mathcal{P}_{\text{map}} \leftarrow \emptyset$
\State $\mathcal{R}'_{\text{all}} \leftarrow \bigcup_{k=1}^n R'_k$
\State Sort $\mathcal{R}'_{\text{all}}$ primarily by $t_{\text{utc}}$ ascending,
       secondarily by source sequence index
\State $\mathcal{B}_{\text{clusters}} \leftarrow
       \textsc{GroupIntoReconciliationClusters}(
         \mathcal{R}'_{\text{all}}, \sim_{\text{dedup}}, \Delta t)$
\For{\textbf{each} cluster $C \in \mathcal{B}_{\text{clusters}}$}
    \State $u \leftarrow \textsc{SynthesizeCanonicalMessage}(C)$
    \State $u.\text{provenance} \leftarrow
           \bigcup_{r' \in C} r'.\text{provenance}$
    \State $\mathcal{P}_{\text{map}}[u.\text{id}] \leftarrow u.\text{provenance}$
    \If{$\textsc{HasUnresolvedConflict}(C)$}
        \State $\mathcal{C}_{\text{conflicts}}.\text{append}(
               \textsc{CreateConflictDescriptor}(C))$
    \EndIf
    \State $\mathcal{U}.\text{append}(u)$
\EndFor
\State $\mathcal{C}_{\text{conflicts}}.\text{extend}(
       \textsc{DetectCrossClusterConflicts}(\mathcal{U}))$
\State $\Omega \leftarrow \big\langle
       \text{provenance\_map}: \mathcal{P}_{\text{map}},\;
       \text{quarantine}: \bigcup \mathcal{Q},\;
       \text{conflicts}: \mathcal{C}_{\text{conflicts}},\;
       \text{stats}: \textsc{ComputeStats}(
         \mathcal{U}, \mathcal{R}'_{\text{all}}, \mathcal{Q}) \big\rangle$
\State \Return $(\mathcal{U}, \Omega)$
\end{algorithmic}
\end{algorithm}

Line 2 is where reconstruction becomes possible at all, and it is the first point
in the pipeline at which fragments meet. Everything before it treats each fragment
in isolation.

Line 3 establishes the total order, with the secondary key supplying determinism
as described above.

Line 4 applies the similarity criterion. The procedure name matters: it groups
records into clusters, and it does not partition them into classes.

Lines 6 to 8 build one canonical record per cluster and union the provenance, then
index it in the map that becomes $\Omega_{\text{provenance}}$. A cluster of one is
the ordinary case, a message that appeared in a single fragment, and it passes
through with its single provenance entry intact.

Lines 9 to 11 handle a cluster whose members are grouped by the similarity
criterion yet disagree in some respect the criterion does not cover. Such a
cluster still yields a canonical record, and it also yields a conflict descriptor.
Detection and reporting are separated from resolution, which is the same
discipline the quarantine rule applies in Normalization: the procedure records
what it could not settle instead of settling it silently.

Line 14 runs the cross-cluster pass of Section~\ref{sec:algorithms:conflicts},
which is the one form of conflict the cluster loop cannot see, and appends its
descriptors to the same log.

Line 15 assembles the report. The quarantine sets arrive at this line untouched,
having been carried through from Normalization for exactly this purpose.

\subsection{Conflict detection}
\label{sec:algorithms:conflicts}

Conflicts arise at two scopes, and Algorithm~\ref{alg:carta_synthesis} detects
each at the line that can see it.

Within a cluster, $\textsc{HasUnresolvedConflict}(C)$ at line 9 holds when the
observations in $C$ disagree on a field the similarity criterion does not
constrain, so that the field-wise precedence rule of
Section~\ref{sec:algorithms:similarity} had to choose between incompatible
values rather than between compatible ones. Sender identity, message type,
attachment metadata beyond the $\cong$ used in
Equation~(\ref{eq:dedup}), and timestamps differing by more than the granularity
the cluster was formed under are all such disagreements. A cluster of one never
satisfies the predicate. The name is deliberate: the question is not whether the
cluster was ambiguous at some point in processing, but whether anything about it
remains unsettled once the declared rules have been applied.

Across clusters, two records may describe genuinely different messages, and so
belong to different clusters, while still standing in contradiction.
$\textsc{DetectCrossClusterConflicts}$ at line 14 is the pass that finds these.
Its criterion applies to records falling in the same coarse bucket but
carrying incompatible sender or ordering constraints:

\begin{equation}
  \begin{split}
    \mathrm{Conflict}(u_i, u_j) \iff{}
      &\lfloor t_{\text{utc}}(u_i) \rfloor_{\Delta t_{\text{conflict}}}
        = \lfloor t_{\text{utc}}(u_j) \rfloor_{\Delta t_{\text{conflict}}} \\
      &{}\land\; u_i \not\sim_{\text{dedup}} u_j \\
      &{}\land\; \big( s(u_i) \neq s(u_j)
        \;\lor\; \mathrm{seq}(u_i, u_j) \text{ is unsatisfiable} \big)
  \end{split}
  \label{eq:conflict}
\end{equation}

Here $\lfloor \cdot \rfloor_{\Delta t_{\text{conflict}}}$ truncates a timestamp
to its containing coarse bucket, and $\mathrm{seq}(u_i, u_j)$ denotes the ordering constraints
the two records' source sequence indices impose: unsatisfiable means the
fragments order the pair in opposite directions, so no single total order
respects both. The second conjunct is what keeps this pass disjoint from the
cluster loop --- a pair the similarity criterion already grouped is one message
seen twice, not two messages in contradiction, and it is handled at line 9.

Coarse buckets are used here because export timestamp precision is itself
unreliable, and demanding exact equality before considering a contradiction would
miss the cases most likely to matter. The consequence of flagging is a report
entry. Neither record is dropped, neither is rewritten, and both remain in
$\mathcal{U}$ with their provenance intact, which is what Invariant 1 requires.
Equation~(\ref{eq:conflict}) detects and reports; it resolves nothing, and no
conforming implementation may use it to drop a record.

The pass is a single line of Algorithm~\ref{alg:carta_synthesis} rather than an
algorithm block of its own because it writes to the same conflict log and needs
no input the procedure does not already hold. Placing it after the loop, rather
than inside, is what the scope requires: it compares canonical records against
one another, so it cannot run before every cluster has produced one.
